% Options for packages loaded elsewhere
% Options for packages loaded elsewhere
\PassOptionsToPackage{unicode}{hyperref}
\PassOptionsToPackage{hyphens}{url}
\PassOptionsToPackage{dvipsnames,svgnames,x11names}{xcolor}
%
\documentclass[
  letterpaper,
  DIV=11,
  numbers=noendperiod]{scrartcl}
\usepackage{xcolor}
\usepackage[margin=1in]{geometry}
\usepackage{amsmath,amssymb}
\setcounter{secnumdepth}{-\maxdimen} % remove section numbering
\usepackage{iftex}
\ifPDFTeX
  \usepackage[T1]{fontenc}
  \usepackage[utf8]{inputenc}
  \usepackage{textcomp} % provide euro and other symbols
\else % if luatex or xetex
  \usepackage{unicode-math} % this also loads fontspec
  \defaultfontfeatures{Scale=MatchLowercase}
  \defaultfontfeatures[\rmfamily]{Ligatures=TeX,Scale=1}
\fi
\usepackage{lmodern}
\ifPDFTeX\else
  % xetex/luatex font selection
  \setmainfont[]{Atkinson Hyperlegible}
  \setmathfont[]{STIX Two Math}
\fi
% Use upquote if available, for straight quotes in verbatim environments
\IfFileExists{upquote.sty}{\usepackage{upquote}}{}
\IfFileExists{microtype.sty}{% use microtype if available
  \usepackage[]{microtype}
  \UseMicrotypeSet[protrusion]{basicmath} % disable protrusion for tt fonts
}{}
\makeatletter
\@ifundefined{KOMAClassName}{% if non-KOMA class
  \IfFileExists{parskip.sty}{%
    \usepackage{parskip}
  }{% else
    \setlength{\parindent}{0pt}
    \setlength{\parskip}{6pt plus 2pt minus 1pt}}
}{% if KOMA class
  \KOMAoptions{parskip=half}}
\makeatother
% Make \paragraph and \subparagraph free-standing
\makeatletter
\ifx\paragraph\undefined\else
  \let\oldparagraph\paragraph
  \renewcommand{\paragraph}{
    \@ifstar
      \xxxParagraphStar
      \xxxParagraphNoStar
  }
  \newcommand{\xxxParagraphStar}[1]{\oldparagraph*{#1}\mbox{}}
  \newcommand{\xxxParagraphNoStar}[1]{\oldparagraph{#1}\mbox{}}
\fi
\ifx\subparagraph\undefined\else
  \let\oldsubparagraph\subparagraph
  \renewcommand{\subparagraph}{
    \@ifstar
      \xxxSubParagraphStar
      \xxxSubParagraphNoStar
  }
  \newcommand{\xxxSubParagraphStar}[1]{\oldsubparagraph*{#1}\mbox{}}
  \newcommand{\xxxSubParagraphNoStar}[1]{\oldsubparagraph{#1}\mbox{}}
\fi
\makeatother


\usepackage{longtable,booktabs,array}
\usepackage{calc} % for calculating minipage widths
% Correct order of tables after \paragraph or \subparagraph
\usepackage{etoolbox}
\makeatletter
\patchcmd\longtable{\par}{\if@noskipsec\mbox{}\fi\par}{}{}
\makeatother
% Allow footnotes in longtable head/foot
\IfFileExists{footnotehyper.sty}{\usepackage{footnotehyper}}{\usepackage{footnote}}
\makesavenoteenv{longtable}
\usepackage{graphicx}
\makeatletter
\newsavebox\pandoc@box
\newcommand*\pandocbounded[1]{% scales image to fit in text height/width
  \sbox\pandoc@box{#1}%
  \Gscale@div\@tempa{\textheight}{\dimexpr\ht\pandoc@box+\dp\pandoc@box\relax}%
  \Gscale@div\@tempb{\linewidth}{\wd\pandoc@box}%
  \ifdim\@tempb\p@<\@tempa\p@\let\@tempa\@tempb\fi% select the smaller of both
  \ifdim\@tempa\p@<\p@\scalebox{\@tempa}{\usebox\pandoc@box}%
  \else\usebox{\pandoc@box}%
  \fi%
}
% Set default figure placement to htbp
\def\fps@figure{htbp}
\makeatother





\setlength{\emergencystretch}{3em} % prevent overfull lines

\providecommand{\tightlist}{%
  \setlength{\itemsep}{0pt}\setlength{\parskip}{0pt}}



 


\linespread{1.05}
\RedeclareSectionCommand[
  afterskip=0.3ex,
  beforeskip=0.8ex
  ]{subsection}
\RedeclareSectionCommand[
  afterskip=0.1ex,
  beforeskip=0.5ex
  ]{subsubsection}
\KOMAoption{captions}{tableheading}
\makeatletter
\@ifpackageloaded{caption}{}{\usepackage{caption}}
\AtBeginDocument{%
\ifdefined\contentsname
  \renewcommand*\contentsname{Table of contents}
\else
  \newcommand\contentsname{Table of contents}
\fi
\ifdefined\listfigurename
  \renewcommand*\listfigurename{List of Figures}
\else
  \newcommand\listfigurename{List of Figures}
\fi
\ifdefined\listtablename
  \renewcommand*\listtablename{List of Tables}
\else
  \newcommand\listtablename{List of Tables}
\fi
\ifdefined\figurename
  \renewcommand*\figurename{Figure}
\else
  \newcommand\figurename{Figure}
\fi
\ifdefined\tablename
  \renewcommand*\tablename{Table}
\else
  \newcommand\tablename{Table}
\fi
}
\@ifpackageloaded{float}{}{\usepackage{float}}
\floatstyle{ruled}
\@ifundefined{c@chapter}{\newfloat{codelisting}{h}{lop}}{\newfloat{codelisting}{h}{lop}[chapter]}
\floatname{codelisting}{Listing}
\newcommand*\listoflistings{\listof{codelisting}{List of Listings}}
\makeatother
\makeatletter
\makeatother
\makeatletter
\@ifpackageloaded{caption}{}{\usepackage{caption}}
\@ifpackageloaded{subcaption}{}{\usepackage{subcaption}}
\makeatother
\usepackage{bookmark}
\IfFileExists{xurl.sty}{\usepackage{xurl}}{} % add URL line breaks if available
\urlstyle{same}
\hypersetup{
  pdftitle={PHIL 342: Introduction to Groups and Choices},
  pdfauthor={Brian Weatherson},
  colorlinks=true,
  linkcolor={blue},
  filecolor={Maroon},
  citecolor={Blue},
  urlcolor={Blue},
  pdfcreator={LaTeX via pandoc}}


\title{PHIL 342: Introduction to Groups and Choices}
\author{Brian Weatherson}
\date{Winter 2026}
\begin{document}
\maketitle


\textbf{Lead Instructor}: Brian Weatherson\\
\textbf{Email}: weath@umich.edu\\
\textbf{Web}: canvas.umich.edu\\
\textbf{Office Hours}: Thursday 9.30-11.30 (Angell Hall, Rm 2207)\\
\textbf{Discussion Section Leader}: Mitch Barrington\\
\textbf{Email}: umitch@umich.edu\\
\textbf{Office Hours}: Monday 1-3 (Angell Hall, Rm 2253)

~

\section{Key Points}\label{key-points}

\begin{itemize}
\tightlist
\item
  This course examines philosophical questions that arise when
  individual choices interact.
\item
  The course has two main parts: pooling (how we form collective
  attitudes from individual attitudes) and modeling (what simple models
  reveal about social structures).
\item
  We'll study voting methods, preference aggregation, opinion pooling,
  game theory, signaling, and models of scientific communities.
\item
  Three textbooks are available electronically through the library:
  \emph{Opinion Pooling} by Lee Elkin and Richard Pettigrew,
  \emph{Modeling Scientific Communities} by Cailin O'Connor, and
  \emph{Essentials of Game Theory} by Kevin Leyton-Brown and Yoav
  Shoham.
\item
  Assessment consists of four short assignments (10\% each), two essays
  (25\% each), and discussion section participation (10\%)
\end{itemize}

\section{Course Description}\label{course-description}

This course introduces topics at the intersection of philosophy, game
theory, and social choice theory. We examine how individual choices and
judgments combine to form group decisions, and what mathematical models
can teach us about social phenomena.

The first half focuses on pooling problems: when we try to aggregate
individual beliefs or preferences into collective judgments, we
encounter surprising paradoxes and impossibility results. We'll explore
voting methods, Arrow's theorem, judgment aggregation, and methods for
pooling probabilistic opinions.

The second half looks at the use mathematical and computational models
to understand social dynamics. The topics we'll cover include focal
points, signaling in markets, iterated games, and agent-based models of
scientific communities and knowledge production.

\section{Required Materials}\label{required-materials}

There are three textbooks for the course. All of them are available
electronically through the library (and the first is open access, so
doesn't even need a library login).

\begin{itemize}
\tightlist
\item
  Lee Elkin and Richard Pettigrew,
  \href{https://www.cambridge.org/core/elements/opinion-pooling/40647AA21F276CF80F67398058682294}{Opinion
  Pooling}
\item
  Cailin O'Connor,
  \href{https://www.cambridge.org/core/elements/abs/modelling-scientific-communities/1ED3515216067E40A37A72094EE3CB15}{Modeling
  Scientific Communities} (if you are off-campus you can access this via
  \href{https://search.lib.umich.edu/catalog/record/99187887203306381}{the
  library})
\item
  Kevin Leyton-Brown and Yoav Shoham,
  \href{https://link-springer-com.proxy.lib.umich.edu/book/10.1007/978-3-031-01545-8}{Essentials
  of Game Theory}
\end{itemize}

\section{Course Requirements}\label{course-requirements}

\begin{itemize}
\tightlist
\item
  There will be 4 short assignments, each worth 10\% of the grade
\item
  There will be 2 essays, each worth 25\% of the grade.
\item
  Lecture participation will not be graded, and we won't take attendance
  here. But 10\% of the grade is for discussion section participation
  (and any other assignments the discussion section leader chooses to
  include).
\end{itemize}

\section{Summary of Grading System}\label{summary-of-grading-system}

\begin{enumerate}
\def\labelenumi{\arabic{enumi}.}
\tightlist
\item
  Short assignments - 4 assignments, 10\% each: 40\%
\item
  Essays - 2 essays, 25\% each: 50\%
\item
  Discussion section participation. (10\%)
\end{enumerate}

\section{Canvas}\label{canvas}

There is a Canvas site for this course, which can be accessed from
\url{https://canvas.umich.edu}. Course documents will be available from
this site. Please make sure that you can access this site. Consult the
site regularly for announcements, including changes to the course
schedule. And there are many tools on the site to communicate with each
other, and with me.

There is also a github site for the course which contains more
documents, and it is at
\url{https://bweatherson.github.io/W26-Phil342-site/}.

\pagebreak

\section{Class Schedule}\label{class-schedule}

Readings should be done before the scheduled class.

\subsection{Week 1: Introduction}\label{week-1-introduction}

\subsubsection{Wednesday, January 07}\label{wednesday-january-07}

\textbf{Topic}: Introduction\\
\textbf{Reading}: None.

\textbf{NOTE}: No discussion section this week.

\subsection{Week 2: Voting}\label{week-2-voting}

\subsubsection{Monday, January 12}\label{monday-january-12}

\textbf{Topic}: Voting Methods\\
\textbf{Reading}: Eric Pacuit,
\href{https://plato.stanford.edu/entries/voting-methods/}{Voting
Methods}

\subsubsection{Wednesday, January 14}\label{wednesday-january-14}

\textbf{Topic}: Arrow's Theorem\\
\textbf{Reading}: Michael Morreau,
\href{https://plato.stanford.edu/archives/win2022/entries/arrows-theorem/}{Arrow's
Theorem}

\subsection{Week 3: Arrow's Theorem}\label{week-3-arrows-theorem}

\subsubsection{Monday, January 19}\label{monday-january-19}

No class for MLK Day

\subsubsection{Wednesday, January 21}\label{wednesday-january-21}

\textbf{Topic}: Responses to Arrow\\
\textbf{Reading}: No new reading

\subsection{Week 4: Sen's Paradox}\label{week-4-sens-paradox}

\subsubsection{Monday, January 26}\label{monday-january-26}

\textbf{Topic}: Sen's Paradox\\
\textbf{Reading}: Amartya Sen,
\href{https://www.jstor.org/stable/1829633?seq=1}{The Impossibility of a
Paretian Liberal}\\
\textbf{Suggested Reading}: Amartya Sen,
\href{https://www.jstor.org/stable/2026284?seq=1}{Liberty and Social
Choice}

\subsubsection{Wednesday, January 28}\label{wednesday-january-28}

No lecture, because I'm away at a conference. There will be a discussion
section.

\textbf{Assignment}: First short assignment due Thursday, January 29 at
5pm.

\subsection{Week 5: Judgment
Aggregation}\label{week-5-judgment-aggregation}

\subsubsection{Monday, February 02}\label{monday-february-02}

\textbf{Topic}: From Belief to Probability\\
\textbf{Reading}: \emph{Opinion Pooling}, Ch. 1 and sections 2.1-2.2.\\
\textbf{Suggested Reading}:

\begin{itemize}
\tightlist
\item
  David Christensen,
  \href{https://philpapers.org/go.pl?id=CHREOD&proxyId=&u=https\%3A\%2F\%2Fphilpapers.org\%2Farchive\%2FCHREOD.pdf}{Epistemology
  of Disagreement: The Good News}
\item
  Adam Elga,
  \href{https://philsci-archive.pitt.edu/2940/1/refdis.pdf}{Reflection
  and Disagreement}
\item
  Thomas Kelly,
  \href{https://philpapers.org/go.pl?id=KELTES-3&proxyId=&u=https\%3A\%2F\%2Fwww.princeton.edu\%2F~tkelly\%2Fesod.pdf}{The
  Epistemic Significance of Disagreement}
\end{itemize}

\subsubsection{Wednesday, February 04}\label{wednesday-february-04}

\textbf{Topic}: Reasons for Aggregation\\
\textbf{Reading}: \emph{Opinion Pooling}, Rest of chapter 2.\\
\textbf{Suggested Reading}:

\begin{itemize}
\tightlist
\item
  Roger White,
  \href{https://philpapers.org/go.pl?id=WHIEP&proxyId=&u=http\%3A\%2F\%2Fweb.mit.edu\%2Frog\%2Fwww\%2Fpapers\%2FEP.pdf}{Epistemic
  Permissiveness}
\item
  Miriam Schoenfield,
  \href{https://philpapers.org/go.pl?id=SCHPTB&proxyId=&u=https\%3A\%2F\%2Fphilpapers.org\%2Farchive\%2FSCHPTB.pdf}{Permission
  to Believe: Why Permissivism Is True and What It Tells Us About
  Irrelevant Influences on Belief}
\end{itemize}

\subsection{Week 6: Constraints on
Aggregation}\label{week-6-constraints-on-aggregation}

\subsubsection{Monday, February 09}\label{monday-february-09}

\textbf{Topic}: Motivating the Constraints\\
\textbf{Reading}: \emph{Opinion Pooling}, Chapter 3.

\subsubsection{Wednesday, February 11}\label{wednesday-february-11}

\textbf{Topic}: Linear and Geometric Pooling \textbf{Reading}:
\emph{Opinion Pooling}, sections 4.1-4.2.\\
\textbf{Suggested Reading}:

\begin{itemize}
\tightlist
\item
  John Hawthorne, Lara Buchak, and Jeffrey Sanford Russell,
  \href{https://philpapers.org/go.pl?id=RUSG&proxyId=&u=https\%3A\%2F\%2Fphilpapers.org\%2Farchive\%2FRUSG.pdf}{Groupthink}
\item
  Jan-Willem Romeijn,
  \href{https://philpapers.org/go.pl?id=ROMAIO&proxyId=&u=https\%3A\%2F\%2Fdx.doi.org\%2F10.1017\%2Fepi.2020.45}{An
  Interpretation of Weights in Linear Opinion Pooling}
\item
  Jean Baccelli \& Rush T. Stewart,
  \href{https://philpapers.org/go.pl?id=BACSFG&proxyId=&u=https\%3A\%2F\%2Fphilpapers.org\%2Farchive\%2FBACSFG.pdf}{Support
  for Geometric Pooling}
\end{itemize}

\subsection{Week 7: Evaluating the
Strategies}\label{week-7-evaluating-the-strategies}

\subsubsection{Monday, February 16}\label{monday-february-16}

\textbf{Topic}: Multiplicative Pooling\\
\textbf{Reading}: \emph{Opinion Pooling}, section 4.3

\subsubsection{Wednesday, February 18}\label{wednesday-february-18}

\textbf{Topic}: Group Rationality\\
\textbf{Reading}: \emph{Opinion Pooling}, chapter 5.\\
\textbf{Suggested Reading}:

\begin{itemize}
\tightlist
\item
  James Joyce,
  \href{https://philpapers.org/go.pl?id=JOYANV-2&proxyId=&u=http\%3A\%2F\%2Ffitelson.org\%2Ftopics\%2Fjoyce_1998.pdf}{A
  Nonpragmatic Vindication of Probabilism}
\item
  Rush T. Stewart \& Ignacio Ojea Quintana,
  \href{https://philpapers.org/go.pl?id=STELAP-15&proxyId=&u=https\%3A\%2F\%2Fdx.doi.org\%2F10.1007\%2Fs10670-017-9894-2}{Learning
  and Pooling, Pooling and Learning}
\end{itemize}

\textbf{Assignment}: Second short assignment due Thursday, February 19
at 5pm.

\subsection{Week 8: Schelling}\label{week-8-schelling}

\subsubsection{Monday, February 23}\label{monday-february-23}

\textbf{Topic}: Focal Points\\
\textbf{Reading}: Judith Mehta, Chris Starmer, and Robert Sugden,
\href{https://www.jstor.org/stable/2118074}{The Nature of Salience: An
Experimental Investigation of Pure Coordination Games}\\
\textbf{Suggested Reading}: Robert Sugden,
\href{https://www.jstor.org/stable/2235016}{A Theory of Focal Points}

\subsubsection{Wednesday, February 25}\label{wednesday-february-25}

\textbf{Topic}: Segregation\\
\textbf{Reading}: Vi Hart and Nicky Case,
\href{https://ncase.me/polygons/}{Parable of the Polygons}\\
\textbf{Suggested Reading}: Thomas Schelling,
\href{https://www.stat.berkeley.edu/~aldous/157/Papers/Schelling_Seg_Models.pdf}{Dynamic
Models of Segregation}

\textbf{Essay}: First essay due Thursday, February 26 at 5pm.

\subsection{Week 9: Game Theory}\label{week-9-game-theory}

\subsubsection{Monday, March 09}\label{monday-march-09}

\textbf{Topic}: One-shot games\\
\textbf{Reading}: \emph{Essentials of Game Theory}, Chapters 1 and 2

\subsubsection{Wednesday, March 11}\label{wednesday-march-11}

\textbf{Topic}: Dynamic games\\
\textbf{Reading}: \emph{Essentials of Game Theory}, Chapters 4 and 5

\subsection{Week 10: Axelrod}\label{week-10-axelrod}

\subsubsection{Monday, March 16}\label{monday-march-16}

\textbf{Topic}: Iterated Prisoners' Dilemma\\
\textbf{Reading}: Robert Axelrod and William D. Hamilton,
\href{https://websites.umich.edu/~axe/research/Axelrod\%20and\%20Hamilton\%20EC\%201981.pdf}{The
Evolution of Cooperation}\\
\textbf{Suggested reading}: Robert Axelrod,
\href{https://esp.mit.edu/download/2dd80247-be60-4e73-b49d-91b5b6988ff0/S12850_axelrod_1980.pdf}{Effective
Choice in the Prisoner's Dilemma}

\subsubsection{Wednesday, March 18}\label{wednesday-march-18}

\textbf{Topic}: Modern Approaches to IPD\\
\textbf{Reading}: Alexander J. Stewart and Joshua B. Plotkin,
\href{https://www.pnas.org/doi/10.1073/pnas.1208087109}{Extortion and
cooperation in the Prisoner's Dilemma}\\
\textbf{Suggested Reading}: William H. Press and Freeman J. Dyson,
\href{https://www.pnas.org/doi/abs/10.1073/pnas.1206569109}{Iterated
Prisoner's Dilemma contains strategies that dominate any evolutionary
opponent}\\
\textbf{Assignment}: Third short assignment due Thursday, March 19 at
5pm.

\subsection{Week 11: Signals}\label{week-11-signals}

\subsubsection{Monday, March 23}\label{monday-march-23}

\textbf{Topic}: Signals in car markets\\
\textbf{Reading}: George Akerlof,
\href{https://www.jstor.org/stable/1879431}{The Market for Lemons}

\subsubsection{Wednesday, March 25}\label{wednesday-march-25}

\textbf{Topic}: Signals in employment markets\\
\textbf{Reading}: Michael Spence,
\href{https://www.nobelprize.org/uploads/2018/06/spence-lecture.pdf}{Signaling
in Retrospect and the Informational Structure of Markets}\\
\textbf{Suggested Reading}: Michael Spence,
\href{https://www.jstor.org/stable/1882010}{Job Market Signaling}

\subsection{Week 12: Social Networks and Scientific
Knowledge}\label{week-12-social-networks-and-scientific-knowledge}

\subsubsection{Monday, March 30}\label{monday-march-30}

\textbf{Topic}: Models of evidence sharing\\
\textbf{Reading}:

\begin{itemize}
\tightlist
\item
  \emph{Modeling Scientific Communities}, Chapter 4 through 4.2.1.
\item
  Kevin J. S. Zollman, \href{https://doi.org/10.1086/525605}{The
  Communication Structure of Epistemic Communities}
\end{itemize}

\subsubsection{Wednesday, April 01}\label{wednesday-april-01}

\textbf{Topic}: Evaluating models of evidence sharing\\
\textbf{Reading}: \emph{Modeling Scientific Communities}, rest of
Chapter 4.\\
\textbf{Suggested Reading}: Sarita Rosenstock, Justin Bruner, and Cailin
O'Connor,
\href{https://cailinoconnor.com/wp-content/uploads/2015/03/In_Epistemic_Networks_Is_Less_Really_Mor-FINAL-VERSION.pdf}{In
Epistemic Networks, Is Less Really More?}.

\subsection{Week 13: The Credit
Economy}\label{week-13-the-credit-economy}

\subsubsection{Monday, April 06}\label{monday-april-06}

\textbf{Topic}: The Matthew Effect\\
\textbf{Reading}:

\begin{itemize}
\tightlist
\item
  \emph{Modeling Scientific Communities}, Chapter 2 through 2.6.
  (Section 2.6 is the most important part.)
\item
  Robert K. Merton,
  \href{https://www.science.org/doi/10.1126/science.159.3810.56}{The
  Matthew Effect in Science}
\end{itemize}

\subsubsection{Wednesday, April 08}\label{wednesday-april-08}

\textbf{Topic}: Credit and Bias\\
\textbf{Reading}:

\begin{itemize}
\tightlist
\item
  \emph{Modeling Scientific Communities}, rest of Chapter 2.
\item
  Remco Heesen,
  \href{https://link.springer.com/article/10.1007/s11229-016-1146-5}{Academic
  Superstars: Competent or Lucky?}
\end{itemize}

\textbf{Suggested Reading}: Hannah Rubin,
\href{https://link.springer.com/article/10.1007/s11098-021-01765-3}{Structural
causes of citation gaps}\\
\textbf{Assignment}: Fourth short assignment due Thursday, April 09 at
5pm.

\subsection{Week 14: Epistemic
Landscapes}\label{week-14-epistemic-landscapes}

\subsubsection{Monday, April 13}\label{monday-april-13}

\textbf{Topic}: Landscapes and exploration\\
\textbf{Reading}:

\begin{itemize}
\tightlist
\item
  \emph{Modeling Scientific Communities}, Chapter 5 through 5.1.
\item
  Michael Weisberg and Ryan Muldoon,
  \href{https://www.jstor.org/stable/10.1086/644786?seq=1}{Epistemic
  Landscapes and the Division of Cognitive Labor}
\end{itemize}

\subsubsection{Wednesday, April 15}\label{wednesday-april-15}

\textbf{Topic}: Rugged epistemic landscapes\\
\textbf{Reading}:

\begin{itemize}
\tightlist
\item
  \emph{Modeling Scientific Communities}, rest of chapter 5.
\item
  Johanna Thoma,
  \href{https://eprints.lse.ac.uk/88129/1/Thoma_Division\%20of\%20Labour_Accepted.pdf}{The
  Epistemic Division of Labor Revisited}
\end{itemize}

There will be a revision class on Monday, April 20, with no new
material, to go over any questions people have about the final essay.

\textbf{FINAL ESSAY}: Due Tuesday, April 21 at 5pm.

\newpage

\section{Plagiarism}\label{plagiarism}

Although team-work, and even co-authorship, is encouraged, plagiarism is
strictly prohibited. You are responsible for making sure that none of
your work is plagiarized. Be sure to cite work that you use, both direct
quotations and paraphrased ideas. Any citation method that is tolerably
clear is permitted, but if you'd like a good description of a citation
scheme that works well in philosophy, look at
\url{http://bit.ly/VDhRJ4}.

You are encouraged to discuss the course material, including
assignments, with your classmates, but all written work that you hand in
under your own name must be your own. If work is handed is as the work
of two people, you are affirming that each person did a fair share of
the work. (Note that when you're submitting work on Canvas, you have to
each submit the paper, even if it is co-authored. That way Canvas knows
that everyone has turned in work.)

You should also be familiar with the academic integrity policies of the
College of Literature, Science \& the Arts at the University of
Michigan, which are available here:
\url{http://www.lsa.umich.edu/academicintegrity/}. Violations of these
policies will be reported to the Office of the Assistant Dean for
Student Academic Affairs, and sanctioned with a course grade of F.

\section{Disability}\label{disability}

The University of Michigan abides by the Americans with Disabilities Act
of 1990, Section 504 of the Rehabilitation Act of 1973, and other
applicable federal and state laws that prohibit discrimination on the
basis of disability, which mandate that reasonable accommodations be
provided for qualified students with disabilities.

If you have a disability, and may require some type of instructional
and/or examination accommodation, please contact me early in the
semester. If you have not already done so, you will also need to
register with the Office of Services for Students with Disabilities. The
office is located at G664 Haven Hall.

For more information on disability services at the University of
Michigan, go to \url{http://ssd.umich.edu}.




\end{document}
